\section{CutMix Algorithm}
\label{appendix:algorithm}

\algnewcommand\Input{\item[\textbf{Input:}]}%
\algnewcommand\Output{\item[\textbf{Output:}]}%
\begin{algorithm*}[t]
  \caption{Pseudo-code of CutMix}
  \begin{algorithmic}[1]
    \For {each iteration}
    \State input, target = get\_minibatch(dataset) \Comment{input is N$\times$C$\times$W$\times$H size tensor, target is N$\times$K size tensor.}
        \If{mode $==$ training}
            \State input\_s, target\_s = shuffle\_minibatch(input, target) \Comment{CutMix starts here.}
            \State lambda = Unif(0,1)
            \State r\_x = Unif(0,W)
            \State r\_y = Unif(0,H)
            \State r\_w = Sqrt(1 - lambda)
            \State r\_h = Sqrt(1 - lambda)
            \State x1 = Round(Clip(r\_x - r\_w / 2, min=0))
            \State x2 = Round(Clip(r\_x + r\_w / 2, max=W))
            \State y1 = Round(Clip(r\_y - r\_h / 2, min=0))
            \State y2 = Round(Clip(r\_y + r\_h / 2, min=H))
            \State input[:, :, x1:x2, y1:y2] = input\_s[:, :, x1:x2, y1:y2]
            \State lambda = 1 - (x2-x1)*(y2-y1)/(W*H) \Comment{Adjust lambda to the exact area ratio.}
            \State target = lambda * target + (1 - lambda) * target\_s \Comment{CutMix ends.}
        \EndIf
        \State output = model\_forward(input)
        \State loss = compute\_loss(output, target)
        \State model\_update()
    \EndFor
  \end{algorithmic}
  \label{alg:cutmix_algorithm}
\end{algorithm*}

We present the code-level description of CutMix algorithm in Algorithm~\ref{alg:cutmix_algorithm}.
{N, C, and K denote the size of minibatch, channel size of input image, and the number of classes.}
First, CutMix shuffles the order of the minibatch input and target along the first axis of the tensors.
And the lambda and the cropping region (x1,x2,y1,y2) are sampled. 
Then, we mix the input and input\_s by replacing the cropping region of input to the region of input\_s.
The target label is also mixed by interpolating method. 

Note that CutMix is easy to implement with few lines (from line $4$ to line $15$), so is very practical algorithm giving significant impact on a wide range of tasks.